\section{Others}
\subsection{Bit Hacks}
\begin{minted}{c++}
// --- Basic ops ---
x & -x                          // LSB value
x & (x - 1)                     // unset LSB
x | (x - 1)                     // set trailing 0s to 1
~x & (x + 1)                    // isolate lowest 0-bit
x && !(x & (x - 1))             // is power of two?
x >> i & 1                      // get bit i
x |= (1 << i)                   // set bit i
x &= ~(1 << i)                  // clear bit i
x ^= (1 << i)                   // toggle bit i
(1 << k) - 1                    // mask of k lowest bits
// --- GCC builtins (suffix ll for 64-bit; undef for 0) ---
__builtin_popcount(x)           // # of 1-bits
__builtin_clz(x)                // leading zeros
__builtin_ctz(x)                // trailing zeros
__builtin_parity(x)             // popcount mod 2
// --- MSB / bit width (32-bit) ---
31 - __builtin_clz(x)           // MSB position (0-idx)
32 - __builtin_clz(x)           // bit width
1 << (31 - __builtin_clz(x))    // MSB value
// --- Iterate set bits of x ---
for (int s = x; s; s &= s - 1) {
  int i = __builtin_ctz(s);     // current bit index
}
// --- Iterate submasks of mask m ---
for (int s = m; ; s = (s - 1) & m) {
  if (s == 0) break;
}
// --- Next number with same popcount ---
int next_popcount(int x) {
  int c = x & -x, r = x + c;
  return r | (((r ^ x) >> 2) / c);
}
// --- Round up to power of two ---
uint32_t next_pow2(uint32_t n) {
  n--; n |= n>>1; n |= n>>2; n |= n>>4;
  n |= n>>8; n |= n>>16; return ++n;
}
// 64-bit: add n |= n>>32
// --- Gray code ---
int gray = x ^ (x >> 1); // binary -> Gray
// Gray -> binary: g = gray; for (int m = g>>1; m; m>>=1) g ^= m;
\end{minted}

\subsection{Zobrist Hashing (XOR Hashing)}
Each element gets a random 64-bit value via \texttt{mt19937\_64}. Hash of a set $=$ XOR of its elements' hashes. $O(1)$ add/remove, collision probability negligible ($\approx 2^{-64}$).

\textbf{Applications:}
\begin{itemize}
  \item \textbf{Subarray anagram detection}: Assign random $rnd[v]$ to each distinct value. Prefix XOR $pre[i+1] = pre[i] \oplus rnd[a[i]]$. Two subarrays $[l,r]$ and $[l',r']$ are anagrams iff $pre[r+1] \oplus pre[l] = pre[r'+1] \oplus pre[l']$.
  \item \textbf{Set equality}: Hash two sets by XORing element hashes; equal $\implies$ equal hash.
  \item \textbf{Tree isomorphism}: $hash[u] = rnd[val[u]] \oplus \bigoplus_{v \in children} hash[v]$ (sort child hashes first for unordered trees).
  \item \textbf{Game state hashing}: Chess/Go engines assign random numbers to each piece+square; board hash $=$ XOR of all active pieces (Zobrist origin).
\end{itemize}
\begin{minted}{c++}
mt19937_64 rng(chrono::steady_clock::now()
  .time_since_epoch().count());
vector<uint64_t> rnd(N);
for (auto &x : rnd) x = rng();

// Prefix XOR for subarray hash
vector<uint64_t> pre(n + 1, 0);
for (int i = 0; i < n; i++)
  pre[i + 1] = pre[i] ^ rnd[a[i]];
// xor of range [l, r] = pre[r+1] ^ pre[l]
\end{minted}

\subsection{Range OR, XOR, AND Tricks}
\begin{minted}{c++}
int64_t xor_0ton(int64_t n) {
  if (n % 4 == 0) return n;
  if (n % 4 == 1) return 1;
  if (n % 4 == 2) return n + 1;
  return 0;
}
int64_t range_xor(int64_t l, int64_t r) {
  return xor_0ton(r) ^ xor_0ton(l - 1);
}
int64_t range_and(int64_t l, int64_t r) {
  int b = 0;
  while ((l >> b) ^ (r >> b)) b += 1;
  return (l >> b) << b;
}
int64_t range_or(int64_t l, int64_t r) {
  int b = 0;
  while ((l >> b) ^ (r >> b)) b += 1;
  return r | ((1LL << b) - 1);
}
\end{minted}

\subsection{Inclusion–Exclusion with Bitmasks}
Count elements satisfying \textbf{none} of $n$ properties $P_1,\dots,P_n$:
$$\Big|\bigcap_{i=1}^n \overline{P_i}\Big| = \sum_{mask=0}^{2^n-1} (-1)^{\text{popcount}(mask)} \cdot f(mask)$$
where $f(mask)$ = count of elements satisfying all properties with bit $=1$ in mask.
\begin{minted}{c++}
template<typename F>
ll inclusion_exclusion(int n, F f) {
  ll ans = 0;
  for (int mask = 0; mask < (1 << n); mask++) {
    ll sign = (__builtin_popcount(mask) & 1) ? -1 : 1;
    ans += sign * f(mask);
  }
  return ans;
}
// Count numbers in [1,N] not divisible by any prime in p[0..n-1]
// auto f = [&](int mask) -> ll {
//   ll prod = 1;
//   for (int i = 0; i < n; i++)
//     if (mask >> i & 1) prod *= p[i];
//   return N / prod;
// };
// ll result = inclusion_exclusion(n, f);
\end{minted}

\subsection{Sum of Digits up to N}
Counts total digit sum of all integers $[1, N]$.
Key idea: For each decimal position with factor $f$, count contribution of each digit value (0-9) to that position across all numbers. $O(\log N)$.
\begin{minted}{c++}
uint64_t sumDigitsUpTo(uint64_t n) {
  if (n <= 0) return 0;
  uint64_t sum = 0;
  for (uint64_t f = 1; f <= n; f *= 10) {
    uint64_t hi = n / (f * 10);
    uint64_t cur = (n / f) % 10;
    uint64_t lo = n % f;
    sum += hi * 45LL * f;
    sum += cur * (lo + 1);
    sum += f * (cur * (cur - 1) / 2);
  }
  return sum;
}
\end{minted}
\subsection{Number of Digits up to N}
Total count of digits written out in $[1, N]$ (e.g. for $N=12$: $1,2,\dots,9,1,0,1,1,1,2$ = $15$ digits).
Key: Numbers with $d$ digits span a range; count their contribution per digit-length. $O(\log N)$.
\begin{minted}{c++}
uint64_t get(uint64_t n) {
  if (n <= 0) return 0;
  uint64_t base = 1, digits = 1, now = 0;
  while (base <= n / 10) {
    base *= 10;
    now += (base - base / 10) * digits;
    digits += 1;
  }
  n -= (base - 1);
  now += n * digits;
  return now;
}
\end{minted}
\subsection{Game Theory \texorpdfstring{\&}{&} Nim}
\textbf{Standard Nim}: XOR sum of piles $= 0 \implies$ Second player wins. $\neq 0 \implies$ First player wins.

\textbf{Subtraction Game}: Can only take up to $k$ stones. XOR sum of $(pile_i \bmod (k+1))$.

\textbf{Index-k Nim}: Can remove from up to $k$ piles at once. Write sizes in binary. Sum each column. If all column sums $\equiv 0 \pmod{k+1}$, second player wins.

\textbf{Mis\`{e}re Nim} (Last player to move LOSES): If all piles have size 1, Second wins if $N$ is ODD. Otherwise, play like Standard Nim (XOR sum).

\subsection{Sprague-Grundy}
$g(v) = \text{mex}\{g(u) : u \text{ reachable from } v\}$. $g=0 \implies$ losing, else winning. XOR combines sub-games, mex computes single state.
\begin{minted}{c++}
// DP template (bool[] avoids set overhead)
int g[N]; g[0] = 0;
for (int i = 1; i <= maxN; i++) {
  bool seen[K] = {};
  for (int m : moves) if (i >= m) seen[g[i - m]] = 1;
  int m = 0; while (seen[m]) m++;
  g[i] = m;
}
// Standard Nim: g(n) = n (XOR pile sizes / 1st wins if != 0)
// Subtraction {1..m}: g(n) = n % (m+1)
// Take & Break: g(n) = mex{ g(a) ^ g(b) | a+b=n }
// Staircase Nim: XOR odd-indexed piles only
// Wythoff: losing pairs = (floor(phi*k), floor(phi^2*k))
// Pattern: DP 100 terms for periodicity.
// g >> 1000 => closed form.
// Watch: mex includes 0; g(0)=0; add moves irrelevant;
// count losing via pairs XOR=0
\end{minted}
